\section{Complete Article Census of \emph{Summa theologiae} I, qq.~50--64 and 106--114}
\label{app:article-census}

This census lists all 118~articles of St.~Thomas Aquinas's treatise on the angels in exact canonical order: Prima Pars questions~50--64 (72~articles, on angelic substance, knowledge, will, creation, grace, sin, and punishment) and questions~106--114 (46~articles, on illumination, speech, hierarchy, government, ministry, guardianship, and demonic assault). Each entry states the article's question, \emph{sed contra} authority, determination, proposition grade, and recorded contrary opinion.

\Needspace{14\baselineskip}
\paragraph{I q.~50: The Substance of the Angels Absolutely Considered (Sec.~5)}
{\scriptsize
\setlength{\tabcolsep}{2.5pt}
\renewcommand{\arraystretch}{1.15}
\begin{longtable}{>{\raggedright\arraybackslash}p{0.05\linewidth}>{\raggedright\arraybackslash}p{0.23\linewidth}>{\raggedright\arraybackslash}p{0.31\linewidth}>{\raggedright\arraybackslash}p{0.12\linewidth}>{\raggedright\arraybackslash}p{0.24\linewidth}}
\toprule
\textbf{Art.} & \textbf{Question Posed \& \emph{Sed Contra}} & \textbf{Aquinas's Determination} & \textbf{Grade} & \textbf{Contrary Opinion} \\
\midrule
\endfirsthead
\toprule
\textbf{Art.} & \textbf{Question Posed \& \emph{Sed Contra}} & \textbf{Aquinas's Determination} & \textbf{Grade} & \textbf{Contrary Opinion} \\
\midrule
\endhead
\bottomrule
\endfoot
\textbf{1} & Whether an angel is altogether incorporeal? (\emph{SC:} Psalm 103:4 (Vulgate)) & Affirmative: angels are entirely incorporeal. & \textbf{Church teaching} & Subtle aerial or fiery bodies entertained by some early Fathers (e.g., Basil, Damascene DFO II.3 comparatively; Augustine allows inquiry in De civ. Dei XXI; cf. De malo q. 16 a. 1) \\
\addlinespace[1.5pt]
\textbf{2} & Whether an angel is composed of matter and form? (\emph{SC:} Pseudo-Dionysius, De Divinis Nominibus IV) & Negative: an angel is not composed of matter and form, but is a subsisting immaterial form. & \textbf{Thomist position} & Universal hylomorphism: Bonaventure, In II Sent. d. 3 p. 1 a. 1 q. 1; Breviloquium II.6 (following Avicebron) \\
\addlinespace[1.5pt]
\textbf{3} & Whether the angels exist in any great number? (\emph{SC:} Daniel 7:10) & Affirmative: angels exist in an exceedingly vast multitude, incomparably surpassing all material multitude. & \textbf{Common teaching} & Aristotelian restriction to the number of celestial movers (Metaph. XII) \\
\addlinespace[1.5pt]
\textbf{4} & Whether the angels differ in species? (\emph{SC:} Aristotle (Metaphysics III) and Pseudo-Dionysius (De Coelesti Hierarchia X)) & Affirmative: every angel differs specifically from every other angel; it is impossible for two angels to belong to the same species. & \textbf{Thomist position} & Multiple angels in one species: Bonaventure, In II Sent. d. 3 p. 1 a. 2 q. 1; Paris 1277 arts. 81, 96; Scotus (haecceitas, In II Sent. d. 3 q. 6) \\
\addlinespace[1.5pt]
\textbf{5} & Whether the angels are incorruptible? (\emph{SC:} Pseudo-Dionysius, De Divinis Nominibus IV) & Affirmative: angels are naturally incorruptible. Because corruption occurs only by separation of form from matter, a subsisting form cannot lose the existence (esse) that belongs to form per se (as a circle cannot lose roundness). & \textbf{Church teaching} & None recorded among Catholic schools (immortality depends on divine conservation, Damascene DFO II.3) \\
\end{longtable}
}

\Needspace{14\baselineskip}
\paragraph{I q.~51: The Angels in Comparison with Bodies (Sec.~5)}
{\scriptsize
\setlength{\tabcolsep}{2.5pt}
\renewcommand{\arraystretch}{1.15}
\begin{longtable}{>{\raggedright\arraybackslash}p{0.05\linewidth}>{\raggedright\arraybackslash}p{0.23\linewidth}>{\raggedright\arraybackslash}p{0.31\linewidth}>{\raggedright\arraybackslash}p{0.12\linewidth}>{\raggedright\arraybackslash}p{0.24\linewidth}}
\toprule
\textbf{Art.} & \textbf{Question Posed \& \emph{Sed Contra}} & \textbf{Aquinas's Determination} & \textbf{Grade} & \textbf{Contrary Opinion} \\
\midrule
\endfirsthead
\toprule
\textbf{Art.} & \textbf{Question Posed \& \emph{Sed Contra}} & \textbf{Aquinas's Determination} & \textbf{Grade} & \textbf{Contrary Opinion} \\
\midrule
\endhead
\bottomrule
\endfoot
\textbf{1} & Whether the Angels Have Bodies Naturally United to Them? (\emph{SC:} Pseudo-Dionysius, De Divinis Nominibus IV) & Negative: angels do not have bodies naturally united to them. & \textbf{Church teaching} & Platonist aerial-body thesis discussed and left open by Augustine (De civ. Dei XXI; cf. De malo q. 16 a. 1) \\
\addlinespace[1.5pt]
\textbf{2} & Whether Angels Assume Bodies? (\emph{SC:} Augustine, De Civitate Dei XVI) & Affirmative: angels sometimes assume real external bodies. & \textbf{Common teaching} & None recorded \\
\addlinespace[1.5pt]
\textbf{3} & Whether the Angels Exercise Functions of Life in the Bodies Assumed? (\emph{SC:} Theological argument (from q. 51 a. 2 ad 3)) & Negative as to what is strictly proper to life; affirmative only as to what vital actions share with non-vital actions (e.g., sound, local movement). Following Aristotle (De Somn. Et Vig. I), only a living subject with the vital faculty can perform genuine vital operations. & \textbf{Common teaching} & None recorded \\
\end{longtable}
}

\Needspace{14\baselineskip}
\paragraph{I q.~52: The Angels in Relation to Place (Sec.~5)}
{\scriptsize
\setlength{\tabcolsep}{2.5pt}
\renewcommand{\arraystretch}{1.15}
\begin{longtable}{>{\raggedright\arraybackslash}p{0.05\linewidth}>{\raggedright\arraybackslash}p{0.23\linewidth}>{\raggedright\arraybackslash}p{0.31\linewidth}>{\raggedright\arraybackslash}p{0.12\linewidth}>{\raggedright\arraybackslash}p{0.24\linewidth}}
\toprule
\textbf{Art.} & \textbf{Question Posed \& \emph{Sed Contra}} & \textbf{Aquinas's Determination} & \textbf{Grade} & \textbf{Contrary Opinion} \\
\midrule
\endfirsthead
\toprule
\textbf{Art.} & \textbf{Question Posed \& \emph{Sed Contra}} & \textbf{Aquinas's Determination} & \textbf{Grade} & \textbf{Contrary Opinion} \\
\midrule
\endhead
\bottomrule
\endfoot
\textbf{1} & Whether an Angel Is in a Place? (\emph{SC:} Liturgy --- Dominican Breviary, Compline Collect) & Affirmative, though equivocally compared with bodies: an angel is in a corporeal place not by dimensive contact (being measured or contained by place), but by the application of his virtual power (contact of power) to a place---thereby containing the place virtually rather than being contained by it. & \textbf{Common teaching} & Scotus distinguishes definitive presence from operative contact (In II Sent. d. 2 q. 6; Paris 1277 art. 25) \\
\addlinespace[1.5pt]
\textbf{2} & Whether an Angel Can Be in Several Places at Once? (\emph{SC:} John Damascene, De Fide Orthodoxa II) & Negative: an angel cannot be in several places simultaneously, because angelic power is finite and extends to one determinate object/place at a time (unlike God's infinite universal power, which is everywhere). & \textbf{Common teaching} & None recorded \\
\addlinespace[1.5pt]
\textbf{3} & Whether Several Angels Can Be at the Same Time in the Same Place? (\emph{SC:} Theological argument (analogy of human souls)) & Negative: two angels cannot occupy the same place simultaneously. & \textbf{Thomist position} & Scotus and later Franciscan/Suarezian authors permit co-location by definitive presence \\
\end{longtable}
}

\Needspace{14\baselineskip}
\paragraph{I q.~53: The Local Movement of the Angels (Sec.~5)}
{\scriptsize
\setlength{\tabcolsep}{2.5pt}
\renewcommand{\arraystretch}{1.15}
\begin{longtable}{>{\raggedright\arraybackslash}p{0.05\linewidth}>{\raggedright\arraybackslash}p{0.23\linewidth}>{\raggedright\arraybackslash}p{0.31\linewidth}>{\raggedright\arraybackslash}p{0.12\linewidth}>{\raggedright\arraybackslash}p{0.24\linewidth}}
\toprule
\textbf{Art.} & \textbf{Question Posed \& \emph{Sed Contra}} & \textbf{Aquinas's Determination} & \textbf{Grade} & \textbf{Contrary Opinion} \\
\midrule
\endfirsthead
\toprule
\textbf{Art.} & \textbf{Question Posed \& \emph{Sed Contra}} & \textbf{Aquinas's Determination} & \textbf{Grade} & \textbf{Contrary Opinion} \\
\midrule
\endhead
\bottomrule
\endfoot
\textbf{1} & Whether an Angel Can Be Moved Locally? (\emph{SC:} Article of faith (Christ's descent into Hell)) & Affirmative: an angel (even a beatified angel) can move locally, though equivocally compared to bodily locomotion. & \textbf{Common teaching} & None recorded \\
\addlinespace[1.5pt]
\textbf{2} & Whether an Angel Passes Through Intermediate Space? (\emph{SC:} Philosophical argument (from the nature of change)) & Distinction: if an angel moves continuously, he must pass through intermediate space; if his movement is non-continuous, he can pass from one extreme to another without traversing the intermediate space. & \textbf{Thomist position} & Some scholastics require traversal of every intermediate space \\
\addlinespace[1.5pt]
\textbf{3} & Whether the Movement of an Angel Is Instantaneous? (\emph{SC:} Philosophical argument (from before and after in change)) & Negative: an angel's local movement is not instantaneous, but takes place in time---continuous time if the motion is continuous, or discrete (non-continuous) time (a succession of angelic `nows') if the motion is discontinuous. & \textbf{Thomist position} & Suarezian/Scotist accounts of continuous duration in local motion \\
\end{longtable}
}

\Needspace{14\baselineskip}
\paragraph{I q.~54: The Knowledge of the Angels (Powers) (Sec.~6)}
{\scriptsize
\setlength{\tabcolsep}{2.5pt}
\renewcommand{\arraystretch}{1.15}
\begin{longtable}{>{\raggedright\arraybackslash}p{0.05\linewidth}>{\raggedright\arraybackslash}p{0.23\linewidth}>{\raggedright\arraybackslash}p{0.31\linewidth}>{\raggedright\arraybackslash}p{0.12\linewidth}>{\raggedright\arraybackslash}p{0.24\linewidth}}
\toprule
\textbf{Art.} & \textbf{Question Posed \& \emph{Sed Contra}} & \textbf{Aquinas's Determination} & \textbf{Grade} & \textbf{Contrary Opinion} \\
\midrule
\endfirsthead
\toprule
\textbf{Art.} & \textbf{Question Posed \& \emph{Sed Contra}} & \textbf{Aquinas's Determination} & \textbf{Grade} & \textbf{Contrary Opinion} \\
\midrule
\endhead
\bottomrule
\endfoot
\textbf{1} & Whether an Angel's Act of Understanding Is His Substance? (\emph{SC:} Metaphysical argument (from I q. 3 a. 4)) & Negative: an angel's act of understanding cannot be its substance. Action is the actuality of an operative power, just as esse is the actuality of essence; only Pure Act (God) is identical with His own existence and operation. & \textbf{Common teaching} & None recorded \\
\addlinespace[1.5pt]
\textbf{2} & Whether in the Angel to Understand Is to Exist? (\emph{SC:} Pseudo-Dionysius, De Divinis Nominibus IV) & Negative: in no angel or creature is understanding identical with existing (esse). & \textbf{Common teaching} & None recorded \\
\addlinespace[1.5pt]
\textbf{3} & Whether an Angel's Power of Intelligence Is His Essence? (\emph{SC:} Pseudo-Dionysius, De Coelesti Hierarchia XI) & Negative: neither in an angel nor in any creature is the operative power (intellect) identical with its essence. & \textbf{Thomist position} & Franciscan identity or formal distinction of soul/angel and its faculties (Bonaventure, Scotus) \\
\addlinespace[1.5pt]
\textbf{4} & Whether There Is an Active and a Passive Intellect in an Angel? (\emph{SC:} Aristotle, De Anima III (argument from phantasms)) & Negative (except equivocally): angels have neither a passive (possible) intellect nor an active (agent) intellect. & \textbf{Thomist position} & None recorded \\
\addlinespace[1.5pt]
\textbf{5} & Whether There Is Only Intellectual Knowledge in the Angels? (\emph{SC:} Gregory the Great, Homilia XXIX in Evangelia (on the Ascension)) & Affirmative: angels possess exclusively intellectual knowledge (only intellect and will, no sensitive or imaginative powers), because sensitive powers operate through bodily organs and angels have no bodies naturally united to them. & \textbf{Common teaching} & None recorded \\
\end{longtable}
}

\Needspace{14\baselineskip}
\paragraph{I q.~55: The Medium of the Angelic Knowledge (Sec.~6)}
{\scriptsize
\setlength{\tabcolsep}{2.5pt}
\renewcommand{\arraystretch}{1.15}
\begin{longtable}{>{\raggedright\arraybackslash}p{0.05\linewidth}>{\raggedright\arraybackslash}p{0.23\linewidth}>{\raggedright\arraybackslash}p{0.31\linewidth}>{\raggedright\arraybackslash}p{0.12\linewidth}>{\raggedright\arraybackslash}p{0.24\linewidth}}
\toprule
\textbf{Art.} & \textbf{Question Posed \& \emph{Sed Contra}} & \textbf{Aquinas's Determination} & \textbf{Grade} & \textbf{Contrary Opinion} \\
\midrule
\endfirsthead
\toprule
\textbf{Art.} & \textbf{Question Posed \& \emph{Sed Contra}} & \textbf{Aquinas's Determination} & \textbf{Grade} & \textbf{Contrary Opinion} \\
\midrule
\endhead
\bottomrule
\endfoot
\textbf{1} & Whether the Angels Know All Things by Their Substance? (\emph{SC:} Pseudo-Dionysius, De Divinis Nominibus IV) & Negative: God alone knows all things by His own infinite essence. & \textbf{Common teaching} & None recorded \\
\addlinespace[1.5pt]
\textbf{2} & Whether the Angels Understand by Species Drawn from Things? (\emph{SC:} Pseudo-Dionysius, De Divinis Nominibus VII) & Negative: angels do not draw intelligible species from sensible things, but receive connatural (innate/concreated) species directly from God together with their intellectual nature. & \textbf{Thomist position} & Scotus admits species acquired from created objects alongside innate species \\
\addlinespace[1.5pt]
\textbf{3} & Whether the Higher Angels Understand by More Universal Species Than the Lower Angels? (\emph{SC:} Pseudo-Dionysius (De Coelesti Hierarchia XII) and Liber de Causis) & Affirmative: higher angels understand by fewer and more universal intelligible species than lower angels. & \textbf{Thomist position} & None recorded \\
\end{longtable}
}

\Needspace{14\baselineskip}
\paragraph{I q.~56: The Angels' Knowledge of Immaterial Things (Sec.~6)}
{\scriptsize
\setlength{\tabcolsep}{2.5pt}
\renewcommand{\arraystretch}{1.15}
\begin{longtable}{>{\raggedright\arraybackslash}p{0.05\linewidth}>{\raggedright\arraybackslash}p{0.23\linewidth}>{\raggedright\arraybackslash}p{0.31\linewidth}>{\raggedright\arraybackslash}p{0.12\linewidth}>{\raggedright\arraybackslash}p{0.24\linewidth}}
\toprule
\textbf{Art.} & \textbf{Question Posed \& \emph{Sed Contra}} & \textbf{Aquinas's Determination} & \textbf{Grade} & \textbf{Contrary Opinion} \\
\midrule
\endfirsthead
\toprule
\textbf{Art.} & \textbf{Question Posed \& \emph{Sed Contra}} & \textbf{Aquinas's Determination} & \textbf{Grade} & \textbf{Contrary Opinion} \\
\midrule
\endhead
\bottomrule
\endfoot
\textbf{1} & Whether an Angel Knows Himself? (\emph{SC:} Augustine, De Genesi ad litteram II) & Affirmative: an angel knows himself directly by his own substance/form, without needing an added species. Because an angel is an immaterial subsisting form, he is actually intelligible in himself and immediately united to his own intellect. & \textbf{Common teaching} & None recorded \\
\addlinespace[1.5pt]
\textbf{2} & Whether One Angel Knows Another? (\emph{SC:} Liber de Causis) & Affirmative: one angel knows another angel---not by its own essence (which is distinct), nor by the other's essence (which only the Trinity penetrates), but by innate intelligible species impressed by the Word of God at creation. & \textbf{Thomist position} & Averroes and some scholastics proposed direct intuition of another angel's essence \\
\addlinespace[1.5pt]
\textbf{3} & Whether an Angel Knows God by His Own Natural Principles? (\emph{SC:} Romans 1:19 (a fortiori from human natural knowledge)) & Affirmative, in a specular/mediate way: an angel naturally knows God not by seeing God's essence directly (which exceeds all created nature, citing I q. 12 a. 4), nor through sensible creatures, but through his own angelic essence as a mirror bearing the impressed image of God. & \textbf{Common teaching} & De ver. q. 8 a. 3 ad 17 treats it as intermediate via impressed species rather than specular \\
\end{longtable}
}

\Needspace{14\baselineskip}
\paragraph{I q.~57: The Angels' Knowledge of Material Things (Sec.~6)}
{\scriptsize
\setlength{\tabcolsep}{2.5pt}
\renewcommand{\arraystretch}{1.15}
\begin{longtable}{>{\raggedright\arraybackslash}p{0.05\linewidth}>{\raggedright\arraybackslash}p{0.23\linewidth}>{\raggedright\arraybackslash}p{0.31\linewidth}>{\raggedright\arraybackslash}p{0.12\linewidth}>{\raggedright\arraybackslash}p{0.24\linewidth}}
\toprule
\textbf{Art.} & \textbf{Question Posed \& \emph{Sed Contra}} & \textbf{Aquinas's Determination} & \textbf{Grade} & \textbf{Contrary Opinion} \\
\midrule
\endfirsthead
\toprule
\textbf{Art.} & \textbf{Question Posed \& \emph{Sed Contra}} & \textbf{Aquinas's Determination} & \textbf{Grade} & \textbf{Contrary Opinion} \\
\midrule
\endhead
\bottomrule
\endfoot
\textbf{1} & Whether the Angels Know Material Things? (\emph{SC:} Theological argument (a fortiori from human intellect)) & Affirmative: angels know material things through their connatural intelligible species. & \textbf{Common teaching} & None recorded \\
\addlinespace[1.5pt]
\textbf{2} & Whether an Angel Knows Singulars? (\emph{SC:} Psalm 90:11 (argument from angelic guardianship)) & Affirmative: angels know singular material things as singular---not merely in universal causes (as an astronomer predicts an eclipse)---through their innate universal species, which are likenesses derived from the Divine essence (cause of both form and matter/individuating principles). & \textbf{Common teaching} & Avicenna restricted separate intelligences to universals \\
\addlinespace[1.5pt]
\textbf{3} & Whether Angels Know the Future? (\emph{SC:} Isaias 41:23) & Distinction: angels know future events in their causes---with certainty for necessary effects (e.g., sunrise), conjecturally (far more acutely than human physicians) for effects occurring in the majority of cases, and not at all for contingent/chance events. & \textbf{Common teaching} & None recorded \\
\addlinespace[1.5pt]
\textbf{4} & Whether Angels Know Secret Thoughts? (\emph{SC:} Jeremias 17:9--10) & Distinction: angels (and demons) can infer secret thoughts indirectly from subtle bodily effects and modifications (far better than physicians reading a pulse, citing Augustine De divin. Daemon.), but God alone directly knows the interior thoughts of the mind and affections of the will, which depend strictly on free will. & \textbf{Common teaching} & None recorded \\
\addlinespace[1.5pt]
\textbf{5} & Whether the Angels Know the Mysteries of Grace? (\emph{SC:} Pseudo-Dionysius (De Coelesti Hierarchia VII) and Isaias 63:1) & Distinction: by natural knowledge angels cannot know mysteries of grace at all (since these depend purely on God's free will, 1 Cor 2:11); by beatific vision in the Word they know mysteries of grace progressively and unequally according to God's revelation (1 Cor 2:10), higher angels grasping deeper mysteries and enlightening lower ones. & \textbf{Common teaching} & None recorded \\
\end{longtable}
}

\Needspace{14\baselineskip}
\paragraph{I q.~58: The Mode of Angelic Knowledge (Sec.~6)}
{\scriptsize
\setlength{\tabcolsep}{2.5pt}
\renewcommand{\arraystretch}{1.15}
\begin{longtable}{>{\raggedright\arraybackslash}p{0.05\linewidth}>{\raggedright\arraybackslash}p{0.23\linewidth}>{\raggedright\arraybackslash}p{0.31\linewidth}>{\raggedright\arraybackslash}p{0.12\linewidth}>{\raggedright\arraybackslash}p{0.24\linewidth}}
\toprule
\textbf{Art.} & \textbf{Question Posed \& \emph{Sed Contra}} & \textbf{Aquinas's Determination} & \textbf{Grade} & \textbf{Contrary Opinion} \\
\midrule
\endfirsthead
\toprule
\textbf{Art.} & \textbf{Question Posed \& \emph{Sed Contra}} & \textbf{Aquinas's Determination} & \textbf{Grade} & \textbf{Contrary Opinion} \\
\midrule
\endhead
\bottomrule
\endfoot
\textbf{1} & Whether the Angel's Intellect Is Sometimes in Potentiality, Sometimes in Act? (\emph{SC:} Augustine, De Genesi ad litteram II) & Distinction (using Aristotle's two modes of intellectual potentiality): (1) as lacking habitual species, an angel is never in potentiality regarding natural knowledge, though he can be in potentiality regarding new divine revelations; (2) as possessing species without actually considering them, an angel can be in potentiality regarding natural knowledge, yet never regarding his uninterrupted beatific vision of the Word. & \textbf{Common teaching} & None recorded \\
\addlinespace[1.5pt]
\textbf{2} & Whether an Angel Can Understand Many Things at the Same Time? (\emph{SC:} Augustine, De Genesi ad litteram IV.32) & Distinction: an angel can understand many things simultaneously insofar as they are grasped under one intelligible species, but not as distinct objects under diverse species. & \textbf{Thomist position} & Scotus and Suarez allow simultaneous actuation of distinct species \\
\addlinespace[1.5pt]
\textbf{3} & Whether an Angel's Knowledge Is Discursive? (\emph{SC:} Pseudo-Dionysius, De Divinis Nominibus VII) & Negative: angelic knowledge is not discursive (does not move from known premises to previously unknown conclusions). & \textbf{Common teaching} & None recorded \\
\addlinespace[1.5pt]
\textbf{4} & Whether the Angels Understand by Composing and Dividing? (\emph{SC:} Pseudo-Dionysius (De Divinis Nominibus VII) and Aristotle (De Anima III)) & Negative: angels do not understand by composing and dividing (forming affirmative/negative judgments stepwise). & \textbf{Thomist position} & None recorded \\
\addlinespace[1.5pt]
\textbf{5} & Whether There Can Be Falsehood in the Intellect of an Angel? (\emph{SC:} Aristotle (De Anima III) and Augustine (De Diversis Quaestionibus LXXXIII, q. 32)) & Distinction: no falsehood or deception can exist per se in any angel, nor even accidentally in good angels (who subordinate every judgment to divine ordinance); in demons, however, error occurs accidentally regarding supernatural matters because their perverted wills judge by natural conditions alone. & \textbf{Common teaching} & None recorded \\
\addlinespace[1.5pt]
\textbf{6} & Whether There Is a ``Morning'' and an ``Evening'' Knowledge in the Angels? (\emph{SC:} Augustine (De Genesi ad litteram IV; De Civitate Dei XI)) & Affirmative: adopting Augustine's interpretation of the six days of Genesis as the single day of angelic cognition directed to six classes of creatures, `morning knowledge' is the angels' knowledge of things in their primordial being in the Word, while `evening knowledge' is their knowledge of creatures as they subsist in their own created natures. & \textbf{Common teaching} & None recorded \\
\addlinespace[1.5pt]
\textbf{7} & Whether the Morning and Evening Knowledge Are One? (\emph{SC:} Augustine, De Genesi ad litteram IV.24) & Distinction: evening knowledge (knowing creatures in their proper nature) is exercised through two media; insofar as angels see creatures' proper natural being in the Word, morning and evening knowledge are essentially one, whereas insofar as they know creatures through innate species, the two are distinct. & \textbf{Thomist position} & De ver. q. 8 a. 16 separates them strictly by uncreated vs. created medium \\
\end{longtable}
}

\Needspace{14\baselineskip}
\paragraph{I q.~59: The Will of the Angels (Sec.~7)}
{\scriptsize
\setlength{\tabcolsep}{2.5pt}
\renewcommand{\arraystretch}{1.15}
\begin{longtable}{>{\raggedright\arraybackslash}p{0.05\linewidth}>{\raggedright\arraybackslash}p{0.23\linewidth}>{\raggedright\arraybackslash}p{0.31\linewidth}>{\raggedright\arraybackslash}p{0.12\linewidth}>{\raggedright\arraybackslash}p{0.24\linewidth}}
\toprule
\textbf{Art.} & \textbf{Question Posed \& \emph{Sed Contra}} & \textbf{Aquinas's Determination} & \textbf{Grade} & \textbf{Contrary Opinion} \\
\midrule
\endfirsthead
\toprule
\textbf{Art.} & \textbf{Question Posed \& \emph{Sed Contra}} & \textbf{Aquinas's Determination} & \textbf{Grade} & \textbf{Contrary Opinion} \\
\midrule
\endhead
\bottomrule
\endfoot
\textbf{1} & Whether There Is Will in the Angels? (\emph{SC:} Augustine, De Trinitate X.11--12) & Affirmative: there is necessarily a will in the angels. & \textbf{Church teaching} & None recorded \\
\addlinespace[1.5pt]
\textbf{2} & Whether in the Angels the Will Differs from the Intellect? (\emph{SC:} Theological argument (diversity of objects)) & Affirmative: in angels the will is a distinct faculty, really different from both their essence/nature and their intellect. & \textbf{Thomist position} & Franciscan formal or virtual distinction rather than real distinction of faculties (Bonaventure, Scotus) \\
\addlinespace[1.5pt]
\textbf{3} & Whether There Is Free-Will in the Angels? (\emph{SC:} Theological argument (a fortiori from human dignity)) & Affirmative: angels possess free-will (liberum arbitrium), and in a higher degree of perfection than men. & \textbf{Church teaching} & None recorded \\
\addlinespace[1.5pt]
\textbf{4} & Whether There Is an Irascible and a Concupiscible Appetite in the Angels? (\emph{SC:} Aristotle, De Anima III) & Negative: angels have no irascible or concupiscible appetite, only an undivided intellective appetite (will). & \textbf{Common teaching} & None recorded \\
\end{longtable}
}

\Needspace{14\baselineskip}
\paragraph{I q.~60: The Love or Dilection of the Angels (Sec.~7)}
{\scriptsize
\setlength{\tabcolsep}{2.5pt}
\renewcommand{\arraystretch}{1.15}
\begin{longtable}{>{\raggedright\arraybackslash}p{0.05\linewidth}>{\raggedright\arraybackslash}p{0.23\linewidth}>{\raggedright\arraybackslash}p{0.31\linewidth}>{\raggedright\arraybackslash}p{0.12\linewidth}>{\raggedright\arraybackslash}p{0.24\linewidth}}
\toprule
\textbf{Art.} & \textbf{Question Posed \& \emph{Sed Contra}} & \textbf{Aquinas's Determination} & \textbf{Grade} & \textbf{Contrary Opinion} \\
\midrule
\endfirsthead
\toprule
\textbf{Art.} & \textbf{Question Posed \& \emph{Sed Contra}} & \textbf{Aquinas's Determination} & \textbf{Grade} & \textbf{Contrary Opinion} \\
\midrule
\endhead
\bottomrule
\endfoot
\textbf{1} & Whether There Is Natural Love or Dilection in an Angel? (\emph{SC:} Augustine, De Trinitate X.1--2) & Affirmative: there is natural love (dilectio naturalis) in the angelic will. What is prior is preserved in what is posterior; since nature precedes intellect and every nature has a built-in inclination toward its good, an intellectual nature possesses a natural inclination within its will. & \textbf{Common teaching} & None recorded \\
\addlinespace[1.5pt]
\textbf{2} & Whether There Is Love of Choice in the Angels? (\emph{SC:} Theological argument (from merit and demerit)) & Affirmative: angels have both natural love (of the end, serving as principle) and love of choice (of things ordered to the end). & \textbf{Common teaching} & None recorded \\
\addlinespace[1.5pt]
\textbf{3} & Whether the Angel Loves Himself with Both Natural Love, and Love of Choice? (\emph{SC:} Aristotle, Nicomachean Ethics IX.8) & Affirmative: an angel loves himself with both natural love (naturally desiring his own subsisting good and perfection) and love of choice (freely choosing particular things beneficial to himself). & \textbf{Common teaching} & None recorded \\
\addlinespace[1.5pt]
\textbf{4} & Whether an Angel Loves Another with Natural Love As He Loves Himself? (\emph{SC:} Ecclesiasticus 13:19) & Affirmative: one angel naturally loves another angel as himself insofar as they share a common angelic nature (wishing well to the other as to himself), though not in respects wherein they differ. & \textbf{Common teaching} & None recorded \\
\addlinespace[1.5pt]
\textbf{5} & Whether an angel by natural love loves God more than he loves himself? (\emph{SC:} Theological argument (moral precepts of natural law)) & Affirmative: by natural love both angel and man love God before themselves and more intensely than themselves. & \textbf{Common teaching} & Some later nominalist/Scotist accounts dispute pure natural love of God above self without charity \\
\end{longtable}
}

\Needspace{14\baselineskip}
\paragraph{I q.~61: The Production of the Angels in the Order of Natural Being (Sec.~7)}
{\scriptsize
\setlength{\tabcolsep}{2.5pt}
\renewcommand{\arraystretch}{1.15}
\begin{longtable}{>{\raggedright\arraybackslash}p{0.05\linewidth}>{\raggedright\arraybackslash}p{0.23\linewidth}>{\raggedright\arraybackslash}p{0.31\linewidth}>{\raggedright\arraybackslash}p{0.12\linewidth}>{\raggedright\arraybackslash}p{0.24\linewidth}}
\toprule
\textbf{Art.} & \textbf{Question Posed \& \emph{Sed Contra}} & \textbf{Aquinas's Determination} & \textbf{Grade} & \textbf{Contrary Opinion} \\
\midrule
\endfirsthead
\toprule
\textbf{Art.} & \textbf{Question Posed \& \emph{Sed Contra}} & \textbf{Aquinas's Determination} & \textbf{Grade} & \textbf{Contrary Opinion} \\
\midrule
\endhead
\bottomrule
\endfoot
\textbf{1} & Whether the angels have a cause of their existence? (\emph{SC:} Psalm 148:2, 5) & Affirmative: angels necessarily have a cause of their existence and were created by God. Because God alone is His own existence (I q. 3 a. 4), every other being exists only by participation, and whatever exists by participation must be caused by what exists essentially. & \textbf{Defined} & Dualist and emanationist errors rejected by the councils \\
\addlinespace[1.5pt]
\textbf{2} & Whether the angel was produced by God from eternity? (\emph{SC:} Proverbs 8:22) & Negative: God alone---Father, Son, and Holy Ghost---is from eternity. Catholic Faith holds without doubt (and condemns the contrary as heretical) that God made creatures from nothing, i.e., after they had not been. & \textbf{Defined} & Peripatetic eternity of separate substances \\
\addlinespace[1.5pt]
\textbf{3} & Whether the angels were created before the corporeal world? (\emph{SC:} Genesis 1:1) & Negative as the more probable view, though the contrary is not heretical/erroneous: angels were more probably created simultaneously with the corporeal world, because angels and bodies together form one universe and no part is perfect separated from its whole (Deut 32:4). & \textbf{Disputed} & Eastern Fathers: Gregory of Nazianzus Or. 38.9, Basil Hexaemeron I.5, John of Damascus DFO II.3 (angels created before material cosmos) \\
\addlinespace[1.5pt]
\textbf{4} & Whether the angels were created in the empyrean heaven? (\emph{SC:} Walafrid Strabo (Strabus, Glossa Ordinaria on Gen 1:1)) & Affirmative (by fittingness): although angels do not depend on bodies for existence or creation, it was fitting that they be created in the highest corporeal place (the empyrean heaven) to signify their presidency over all corporeal nature (citing Isidore on Deut 10:14). & \textbf{Thomist position} & Patristic hesitation over localizing immaterial spirits prior to material cosmos \\
\end{longtable}
}

\Needspace{14\baselineskip}
\paragraph{I q.~62: The Perfection of the Angels in the Order of Grace and Glory (Sec.~7)}
{\scriptsize
\setlength{\tabcolsep}{2.5pt}
\renewcommand{\arraystretch}{1.15}
\begin{longtable}{>{\raggedright\arraybackslash}p{0.05\linewidth}>{\raggedright\arraybackslash}p{0.23\linewidth}>{\raggedright\arraybackslash}p{0.31\linewidth}>{\raggedright\arraybackslash}p{0.12\linewidth}>{\raggedright\arraybackslash}p{0.24\linewidth}}
\toprule
\textbf{Art.} & \textbf{Question Posed \& \emph{Sed Contra}} & \textbf{Aquinas's Determination} & \textbf{Grade} & \textbf{Contrary Opinion} \\
\midrule
\endfirsthead
\toprule
\textbf{Art.} & \textbf{Question Posed \& \emph{Sed Contra}} & \textbf{Aquinas's Determination} & \textbf{Grade} & \textbf{Contrary Opinion} \\
\midrule
\endhead
\bottomrule
\endfoot
\textbf{1} & Whether the Angels Were Created in Beatitude? (\emph{SC:} Theological argument (from stability of beatitude and the fall)) & Distinction: angels were created in natural beatitude (ultimate perfection attainable by natural power, possessed instantly without discursive progress, citing Aristotle Ethic. X), but were not created in supernatural beatitude (the beatific vision of God's essence, which exceeds every created nature and is attained as an end rather than possessed from the start). & \textbf{Common teaching} & None recorded \\
\addlinespace[1.5pt]
\textbf{2} & Whether an Angel Needs Grace in Order to Turn to God? (\emph{SC:} Romans 6:23) & Affirmative: angels needed supernatural grace to turn to God as the object of supernatural beatitude (the vision of His essence), because no created will can incline toward an end above its nature unless moved by a supernatural principle (just as fire cannot generate flesh unless used instrumentally by a living soul). & \textbf{Church teaching} & None recorded \\
\addlinespace[1.5pt]
\textbf{3} & Whether the Angels Were Created in Grace? (\emph{SC:} Augustine, De Civitate Dei XII.9) & Affirmative (as the more probable view, noting conflicting theological opinions): it is more probable and more consonant with the Saints that angels were created in sanctifying grace, just as Augustine holds that natural things were created with seminal forms (grace being the `seed' of glory, 1 John 3:9). & \textbf{Disputed} & Peter Lombard (II Sent. d. 3) and Bonaventure (In II Sent. d. 3; Breviloquium II.7) hold creation in natural integrity prior to grace \\
\addlinespace[1.5pt]
\textbf{4} & Whether an Angel Merits His Beatitude? (\emph{SC:} Apocalypse 21:17) & Affirmative: an angel merited his beatitude. & \textbf{Common teaching} & None recorded \\
\addlinespace[1.5pt]
\textbf{5} & Whether the Angel Obtained Beatitude Immediately After One Act of Merit? (\emph{SC:} Luke 20:36 (analogy with separated human souls)) & Affirmative: an angel obtained beatitude immediately after one meritorious act of charity. & \textbf{Thomist position} & Some medieval authors posited a longer interval (mora) of trial \\
\addlinespace[1.5pt]
\textbf{6} & Whether the Angels Receive Grace and Glory According to the Degree of Their Natural Gifts? (\emph{SC:} Peter Lombard, Sentences II d. 3) & Affirmative: it is reasonable to hold that angels received grace and glory in proportion to their natural gifts---both on God's part (divine wisdom crafting nobler stones for the more ornate part of the house) and on the angel's part (lacking conflicting bodily passions, higher angels turned to God with their entire greater natural energy). & \textbf{Thomist position} & Scotus and Suarez deny a necessary proportion between natural rank and degree of glory \\
\addlinespace[1.5pt]
\textbf{7} & Whether Natural Knowledge and Love Remain in the Beatified Angels? (\emph{SC:} Theological argument (grace/glory preserves nature)) & Affirmative: natural knowledge and natural love remain in beatified angels alongside beatific vision and charity, because nature underlies beatitude as prior to posterior, and what is prior is preserved in its perfection. & \textbf{Common teaching} & None recorded \\
\addlinespace[1.5pt]
\textbf{8} & Whether a Beatified Angel Can Sin? (\emph{SC:} Augustine, De Genesi ad litteram XI) & Negative: a beatified angel cannot sin. & \textbf{Church teaching} & Origenist mutability of all rational spirits (condemned Constantinople 543/553) \\
\addlinespace[1.5pt]
\textbf{9} & Whether the Beatified Angels Advance in Beatitude? (\emph{SC:} Theological argument (viatores vs. comprehensores)) & Negative as to essential beatitude: every rational creature is directed by divine predestination to a determinate finite degree of vision of God's infinite essence; once that final term is reached, essential beatitude cannot increase (though accidental joy increases until Judgment Day as men are saved). & \textbf{Common teaching} & None recorded \\
\end{longtable}
}

\Needspace{14\baselineskip}
\paragraph{I q.~63: The Malice of the Angels with Regard to Sin (Sec.~8)}
{\scriptsize
\setlength{\tabcolsep}{2.5pt}
\renewcommand{\arraystretch}{1.15}
\begin{longtable}{>{\raggedright\arraybackslash}p{0.05\linewidth}>{\raggedright\arraybackslash}p{0.23\linewidth}>{\raggedright\arraybackslash}p{0.31\linewidth}>{\raggedright\arraybackslash}p{0.12\linewidth}>{\raggedright\arraybackslash}p{0.24\linewidth}}
\toprule
\textbf{Art.} & \textbf{Question Posed \& \emph{Sed Contra}} & \textbf{Aquinas's Determination} & \textbf{Grade} & \textbf{Contrary Opinion} \\
\midrule
\endfirsthead
\toprule
\textbf{Art.} & \textbf{Question Posed \& \emph{Sed Contra}} & \textbf{Aquinas's Determination} & \textbf{Grade} & \textbf{Contrary Opinion} \\
\midrule
\endhead
\bottomrule
\endfoot
\textbf{1} & Whether the evil of fault can be in the angels? (\emph{SC:} Job 4:18) & Affirmative: every rational creature---angel or human---considered in its own nature is capable of sin (impeccability belongs only to God by nature, and to creatures solely by grace). & \textbf{Defined} & Manichean/Cathar dualism of an evil principle \\
\addlinespace[1.5pt]
\textbf{2} & Whether only the sin of pride and envy can exist in an angel? (\emph{SC:} Augustine, De Civitate Dei XIV.3) & Distinction: by guilt (reatus), demons incur all sins by tempting men to them; by personal affection (affectus), only the spiritual sins of pride (first) and envy (consequently) can exist in them. & \textbf{Common teaching} & Early Watchers tradition of carnal sin (Justin 2 Apol. 5), decisively rejected by Augustine (Civ. Dei XV.23) \\
\addlinespace[1.5pt]
\textbf{3} & Whether the devil desired to be as God? (\emph{SC:} Isaias 14:13--14 and Pseudo-Augustine (Ambrosiaster, De Quaest. Vet. Test. CXIII)) & Distinction: the devil did not desire to be equal to God (knowing naturally this is impossible and would destroy his own species, as an ass does not desire to become a horse), but he did sin by desiring to be like God inordinately---seeking final beatitude by his own natural power (either resting in natural perfection and rejecting grace, or seeking supernatural beatitude by his own power rather than divine gift) and consequently coveting autonomous dominion over others. & \textbf{Common teaching} & Scotus (Ordinatio II d. 6 q. 2) traces the first sin to inordinate self-love (affectio commodi / luxuria spiritualis); Suarez (De angelis VII.13) and later Scotists emphasize refusal of the revealed Incarnation / hypostatic union \\
\addlinespace[1.5pt]
\textbf{4} & Whether any demons are naturally wicked? (\emph{SC:} Pseudo-Dionysius, De Divinis Nominibus IV) & Negative: no demon is naturally wicked. & \textbf{Defined} & Manichean and Cathar dualism \\
\addlinespace[1.5pt]
\textbf{5} & Whether the devil was wicked by the fault of his own will in the first instant of his creation? (\emph{SC:} Genesis 1:31) & Negative: the devil neither was nor could be wicked in the very first instant of his creation. & \textbf{Thomist position} & De malo q. 16 a. 4 shifts the proof from efficient causality to non-simultaneous cognition, excluding instant-1 merit as well; Augustine leaves first-instant sin open in Civ. Dei XI.15 \\
\addlinespace[1.5pt]
\textbf{6} & Whether there was any interval between the creation and the fall of the angel? (\emph{SC:} John 8:44 and Augustine (De Civitate Dei XI.15)) & Negative (as the more probable opinion): there was no extended interval; the devil sinned immediately after the first instant of his creation (in the second discrete angelic instant). & \textbf{Thomist position} & Longer temporal interval (mora) favored by Bonaventure (Breviloquium II.7) and some Fathers \\
\addlinespace[1.5pt]
\textbf{7} & Whether the highest angel among those who sinned was the highest of all? (\emph{SC:} Gregory the Great, Homilia XXXIV in Evangelia) & Affirmative as the more probable view (without condemning Damascene's contrary view as unorthodox): considering proneness to sin, Damascene (De Fide Orth. II) placed the chief fallen angel among the lower angels presiding over earth (agreeing with Platonists in Augustine, De Civ. Dei VII, X); considering motive for pride (excellence), Gregory's view that he was the very highest of all angels is more probable because angelic sin came from free choice, not proneness. & \textbf{Thomist position} & John of Damascus (DFO II.4) holds the fallen chief presided over the terrestrial order, not the highest seraphic rank \\
\addlinespace[1.5pt]
\textbf{8} & Whether the sin of the highest angel was the cause of the others sinning? (\emph{SC:} Apocalypse 12:4) & Affirmative: the sin of the highest fallen angel caused the sin of the other demons---not by compulsion, but by inducement/exhortation. This is confirmed by all demons being subjected to him in punishment (Matt 25:41; 2 Pet 2:19). & \textbf{Common teaching} & None recorded \\
\addlinespace[1.5pt]
\textbf{9} & Whether those who sinned were as many as those who remained firm? (\emph{SC:} 4 Kings 6:16) & Negative (more stood firm than fell): more angels remained steadfast than sinned, because sin is contrary to natural inclination, and what is against nature happens less frequently than what accords with nature. & \textbf{Thomist position} & None recorded \\
\end{longtable}
}

\Needspace{14\baselineskip}
\paragraph{I q.~64: The Punishment of the Demons (Sec.~8)}
{\scriptsize
\setlength{\tabcolsep}{2.5pt}
\renewcommand{\arraystretch}{1.15}
\begin{longtable}{>{\raggedright\arraybackslash}p{0.05\linewidth}>{\raggedright\arraybackslash}p{0.23\linewidth}>{\raggedright\arraybackslash}p{0.31\linewidth}>{\raggedright\arraybackslash}p{0.12\linewidth}>{\raggedright\arraybackslash}p{0.24\linewidth}}
\toprule
\textbf{Art.} & \textbf{Question Posed \& \emph{Sed Contra}} & \textbf{Aquinas's Determination} & \textbf{Grade} & \textbf{Contrary Opinion} \\
\midrule
\endfirsthead
\toprule
\textbf{Art.} & \textbf{Question Posed \& \emph{Sed Contra}} & \textbf{Aquinas's Determination} & \textbf{Grade} & \textbf{Contrary Opinion} \\
\midrule
\endhead
\bottomrule
\endfoot
\textbf{1} & Whether the Demons' Intellect Is Darkened by Privation of the Knowledge of All Truth? (\emph{SC:} Pseudo-Dionysius, De Divinis Nominibus IV) & Negative (with a threefold distinction of knowledge): (1) natural knowledge is neither removed nor diminished in demons (because angelic simplicity prevents stripping natural powers); (2) speculative knowledge of grace (revelation of divine secrets) is diminished, imparted only as needed through angels or temporal effects (Augustine, De Civ. Dei IX.21); (3) affective knowledge of grace (gift of wisdom producing love of God) is utterly lost. & \textbf{Common teaching} & None recorded \\
\addlinespace[1.5pt]
\textbf{2} & Whether the Will of the Demons Is Obstinate in Evil? (\emph{SC:} Psalm 73:23 (Vulgate)) & Affirmative: against Origen's condemned opinion (Peri Archon I.6, universal mutability of created wills), Catholic Faith firmly holds that good angels are confirmed in good and demons are permanently obstinate in evil. & \textbf{Defined} & Origenist apokatastasis / eventual restoration of demons (condemned 543/553); Aquinas notes God's potentia absoluta in De malo q. 16 a. 5 ad 13 \\
\addlinespace[1.5pt]
\textbf{3} & Whether There Is Sorrow in the Demons? (\emph{SC:} Apocalypse 18:7 (a fortiori from human sin)) & Affirmative (as a simple act of the will, not a bodily passion): sorrow is the will's resistance to what is or what is not; demons suffer intense volitional sorrow because they are deprived of naturally desired happiness, curbed in their wicked designs, and envious of those who are saved. & \textbf{Common teaching} & None recorded \\
\addlinespace[1.5pt]
\textbf{4} & Whether Our Atmosphere Is the Demons' Place of Punishment? (\emph{SC:} Augustine, De Genesi ad litteram III.10) & Affirmative (until the Last Judgment): demons have a twofold penal place---hell (due to their sin, where some already torment the damned) and this darksome atmosphere (assigned by Providence until Judgment Day so they may test and exercise men, preserving their utility in the cosmic order); after the Last Judgment all demons and damned will be confined to hell. & \textbf{Common teaching} & None recorded \\
\end{longtable}
}

\Needspace{14\baselineskip}
\paragraph{I q.~106: How One Creature Moves Another (Angelic Illumination) (Sec.~9)}
{\scriptsize
\setlength{\tabcolsep}{2.5pt}
\renewcommand{\arraystretch}{1.15}
\begin{longtable}{>{\raggedright\arraybackslash}p{0.05\linewidth}>{\raggedright\arraybackslash}p{0.23\linewidth}>{\raggedright\arraybackslash}p{0.31\linewidth}>{\raggedright\arraybackslash}p{0.12\linewidth}>{\raggedright\arraybackslash}p{0.24\linewidth}}
\toprule
\textbf{Art.} & \textbf{Question Posed \& \emph{Sed Contra}} & \textbf{Aquinas's Determination} & \textbf{Grade} & \textbf{Contrary Opinion} \\
\midrule
\endfirsthead
\toprule
\textbf{Art.} & \textbf{Question Posed \& \emph{Sed Contra}} & \textbf{Aquinas's Determination} & \textbf{Grade} & \textbf{Contrary Opinion} \\
\midrule
\endhead
\bottomrule
\endfoot
\textbf{1} & Whether one angel enlightens another? (\emph{SC:} Pseudo-Dionysius, De Coelesti Hierarchia VIII) & Affirmative: one angel enlightens another by communicating the manifestation of known truth (Eph 5:13; 3:8). & \textbf{Common teaching} & De ver. q. 9 a. 1 ad 10 emphasizes no new species are infused \\
\addlinespace[1.5pt]
\textbf{2} & Whether one angel moves another angel's will? (\emph{SC:} Theological argument (God alone bestows righteousness)) & Negative as to interior efficacy; affirmative only by persuasion. & \textbf{Common teaching} & None recorded \\
\addlinespace[1.5pt]
\textbf{3} & Whether an inferior angel can enlighten a superior angel? (\emph{SC:} Pseudo-Dionysius (De Coelesti Hierarchia IV; De Ecclesiastica Hierarchia V)) & Negative: inferior angels never enlighten superior angels, nor does God ever bypass the spiritual hierarchy to enlighten a lower angel independently of the higher. & \textbf{Thomist position} & None recorded \\
\addlinespace[1.5pt]
\textbf{4} & Whether the superior angel enlightens the inferior as regards all he himself knows? (\emph{SC:} Gregory the Great / Peter Lombard (Sent. II d. 9; cf. Hom. XXXIV in Ev.) and Pseudo-Dionysius (Coel. Hier. XV)) & Affirmative (communicating all they know as far as the inferior's capacity allows): good is naturally self-diffusive (cf. 1 Pet 4:10), so superior angels generously impart their knowledge to inferiors; yet inferiors receive it less eminently, preserving hierarchical rank (like a pupil learning from a master). & \textbf{Thomist position} & None recorded \\
\end{longtable}
}

\Needspace{14\baselineskip}
\paragraph{I q.~107: The Speech of the Angels (Sec.~9)}
{\scriptsize
\setlength{\tabcolsep}{2.5pt}
\renewcommand{\arraystretch}{1.15}
\begin{longtable}{>{\raggedright\arraybackslash}p{0.05\linewidth}>{\raggedright\arraybackslash}p{0.23\linewidth}>{\raggedright\arraybackslash}p{0.31\linewidth}>{\raggedright\arraybackslash}p{0.12\linewidth}>{\raggedright\arraybackslash}p{0.24\linewidth}}
\toprule
\textbf{Art.} & \textbf{Question Posed \& \emph{Sed Contra}} & \textbf{Aquinas's Determination} & \textbf{Grade} & \textbf{Contrary Opinion} \\
\midrule
\endfirsthead
\toprule
\textbf{Art.} & \textbf{Question Posed \& \emph{Sed Contra}} & \textbf{Aquinas's Determination} & \textbf{Grade} & \textbf{Contrary Opinion} \\
\midrule
\endhead
\bottomrule
\endfoot
\textbf{1} & Whether one angel speaks to another? (\emph{SC:} 1 Corinthians 13:1) & Affirmative: angels speak to one another by interior spiritual speech. & \textbf{Common teaching} & Scotus analyzes speech via impressed intelligible species rather than bare volitional unveiling \\
\addlinespace[1.5pt]
\textbf{2} & Whether the inferior angel speaks to the superior? (\emph{SC:} Pseudo-Dionysius, De Coelesti Hierarchia VII) & Affirmative: inferior angels can speak to superior angels. & \textbf{Thomist position} & De ver. q. 9 a. 5 frames the distinction by object presence vs. intellectual strength \\
\addlinespace[1.5pt]
\textbf{3} & Whether an angel speaks to God? (\emph{SC:} Zacharias 1:12) & Affirmative (in a receptive/reverential sense, not an informative sense): an angel cannot speak to God to impart knowledge or inform Him, but does speak to God by ordering his mind toward God to receive---either consulting the Divine will on new works to be done, or rising in perpetual praise and admiration of the incomprehensible Divine excellence (Gregory, Moralia II). & \textbf{Common teaching} & None recorded \\
\addlinespace[1.5pt]
\textbf{4} & Whether local distance influences the angelic speech? (\emph{SC:} Luke 16:24 (a fortiori from the rich man and Abraham)) & Negative: local distance has zero effect on angelic speech, because angelic speech is a purely intellectual operation abstracting completely from `here' and `now' (without even the accidental phantasm-link of human thought). & \textbf{Common teaching} & None recorded \\
\addlinespace[1.5pt]
\textbf{5} & Whether all the angels know what one speaks to another? (\emph{SC:} Theological argument (a fortiori from human privacy of speech)) & Negative: not all angels know what one angel speaks to another. & \textbf{Common teaching} & Some scholastics debated whether an unveiled concept becomes publicly visible to all \\
\end{longtable}
}

\Needspace{14\baselineskip}
\paragraph{I q.~108: The Angelic Degrees of Hierarchies and Orders (Sec.~9)}
{\scriptsize
\setlength{\tabcolsep}{2.5pt}
\renewcommand{\arraystretch}{1.15}
\begin{longtable}{>{\raggedright\arraybackslash}p{0.05\linewidth}>{\raggedright\arraybackslash}p{0.23\linewidth}>{\raggedright\arraybackslash}p{0.31\linewidth}>{\raggedright\arraybackslash}p{0.12\linewidth}>{\raggedright\arraybackslash}p{0.24\linewidth}}
\toprule
\textbf{Art.} & \textbf{Question Posed \& \emph{Sed Contra}} & \textbf{Aquinas's Determination} & \textbf{Grade} & \textbf{Contrary Opinion} \\
\midrule
\endfirsthead
\toprule
\textbf{Art.} & \textbf{Question Posed \& \emph{Sed Contra}} & \textbf{Aquinas's Determination} & \textbf{Grade} & \textbf{Contrary Opinion} \\
\midrule
\endhead
\bottomrule
\endfoot
\textbf{1} & Whether all the angels are of one hierarchy? (\emph{SC:} Pseudo-Dionysius, De Coelesti Hierarchia VI) & Distinction: viewed on the part of the Prince (God), all rational creatures (angels and men) form one sacred principality (Augustine, De Civ. Dei XII.1); viewed on the part of the governed multitude and their modes of receiving divine illumination, there are three angelic hierarchies (and a distinct human/ecclesiastical hierarchy receiving illumination under sensible veils, Coel. Hier. I). & \textbf{Common teaching} & None recorded \\
\addlinespace[1.5pt]
\textbf{2} & Whether there are several orders in one hierarchy? (\emph{SC:} Ephesians 1:20--21) & Affirmative: each of the three angelic hierarchies necessarily contains three distinct orders (summit, middle, and base---making nine orders total, Dionysius Coel. Hier. VI). & \textbf{Common teaching} & Basil, Chrysostom, and Dionysius (CH VI) leave open whether unrevealed angelic ranks exist beyond Scripture's nine names \\
\addlinespace[1.5pt]
\textbf{3} & Whether there are many angels in one order? (\emph{SC:} Isaias 6:3) & Affirmative (according to our general human classification): because our knowledge of angels is imperfect (Dionysius, Coel. Hier. VI), we group many angels into one order by general similarity of office; yet if we knew their distinctions perfectly, we would see that every single angel (being its own species) has its own proper individual office and rank. & \textbf{Common teaching} & None recorded \\
\addlinespace[1.5pt]
\textbf{4} & Whether the distinction of hierarchies and orders comes from the angelic nature? (\emph{SC:} Peter Lombard, Sentences II d. 9) & Distinction: relative to their natural end (natural knowledge and love of God), angelic orders are distinguished by natural gifts; relative to their supernatural end (beatific vision), they are distinguished adequately by gifts of grace, and dispositively by natural gifts---since angels receive grace in proportion to natural capacity (I q. 62 a. 6), unlike men, whose ecclesiastical orders rest on grace alone. & \textbf{Thomist position} & Peter Lombard (II Sent. d. 9) emphasizes grace/gifts alone \\
\addlinespace[1.5pt]
\textbf{5} & Whether the orders of the angels are properly named? (\emph{SC:} Sacred Scripture (Isa 6:2; Ezek 1 \& 10; Col 1:16; Eph 1:21; Jude 9)) & Affirmative: each order is properly named from the spiritual perfection that belongs to it by way of property (adequate and proportionate to its grade), rather than by excess (higher orders possess lower perfections eminently) or by mere participation (lower orders share higher perfections imperfectly). & \textbf{Common teaching} & Parker 1899 English renders Dynameis as Powers and Exousiai as Authorities \\
\addlinespace[1.5pt]
\textbf{6} & Whether the grades of the orders are properly assigned? (\emph{SC:} Pseudo-Dionysius, De Coelesti Hierarchia VII (and VIII--IX)) & Affirmative: both Dionysius's ordering (grounded in Eph 1:20--21: Principalities--Powers--Virtues--Dominations ascending) and Gregory's ordering (Hom. XXXIV in Ev., grounded in Col 1:16: Thrones--Dominations--Principalities--Powers descending, placing Virtues below Powers) are rational and substantively congruent once their distinct naming perspectives are compared. & \textbf{Disputed} & Pseudo-Dionysius CH VI--IX (Dom./Virt./Pow.; Princ./Arch./Ang.) vs. Gregory Hom. 34.7 (Dom./Princ./Pow.; Virt./Arch./Ang.) and Moralia XXXII.23.48 (Dom./Princ./Pow./Virt. ordering) \\
\addlinespace[1.5pt]
\textbf{7} & Whether the orders will outlast the Day of Judgment? (\emph{SC:} Judges 5:20) & Affirmative as to distinction of grades and beatific offices; negative only as to temporal offices of leading wayfarers to salvation. & \textbf{Common teaching} & None recorded \\
\addlinespace[1.5pt]
\textbf{8} & Whether men are taken up into the angelic orders? (\emph{SC:} Matthew 22:30) & Affirmative as to grade of grace and glory (negative as to grade of nature): human nature never changes into angelic nature, yet because grace is the perfecting principle of angelic orders and is granted to men according to God's liberality rather than natural rank, all saved men---not merely virgins or a select group---are taken up into the angelic orders according to their merited degrees of glory (Luke 20:36; Augustine, De Civ. Dei XII.9). & \textbf{Common teaching} & None recorded \\
\end{longtable}
}

\Needspace{14\baselineskip}
\paragraph{I q.~109: The Ordering of the Bad Angels (Sec.~9)}
{\scriptsize
\setlength{\tabcolsep}{2.5pt}
\renewcommand{\arraystretch}{1.15}
\begin{longtable}{>{\raggedright\arraybackslash}p{0.05\linewidth}>{\raggedright\arraybackslash}p{0.23\linewidth}>{\raggedright\arraybackslash}p{0.31\linewidth}>{\raggedright\arraybackslash}p{0.12\linewidth}>{\raggedright\arraybackslash}p{0.24\linewidth}}
\toprule
\textbf{Art.} & \textbf{Question Posed \& \emph{Sed Contra}} & \textbf{Aquinas's Determination} & \textbf{Grade} & \textbf{Contrary Opinion} \\
\midrule
\endfirsthead
\toprule
\textbf{Art.} & \textbf{Question Posed \& \emph{Sed Contra}} & \textbf{Aquinas's Determination} & \textbf{Grade} & \textbf{Contrary Opinion} \\
\midrule
\endhead
\bottomrule
\endfoot
\textbf{1} & Whether there are orders among the demons? (\emph{SC:} Ephesians 6:12) & Distinction: relative to consummate glory, demons never were in angelic orders; relative to initial habitual grace, they were in them and fell away (I q. 62 a. 3); relative to their grades of nature, demons still remain in their orders, since they retain their natural gifts intact (Dionysius, Div. Nom. IV). & \textbf{Common teaching} & None recorded \\
\addlinespace[1.5pt]
\textbf{2} & Whether among the demons there is precedence? (\emph{SC:} Gloss on 1 Corinthians 15:24) & Affirmative: there is genuine precedence (praelatio) among demons grounded in their natural inequality of rank (action follows nature, so subordinate natures have subordinate actions) and sustained by Divine wisdom, which leaves nothing inordinate (Wis 8:1). & \textbf{Common teaching} & None recorded \\
\addlinespace[1.5pt]
\textbf{3} & Whether there is enlightenment in the demons? (\emph{SC:} Ecclesiasticus 34:4 (linked to I q. 106 a. 1)) & Negative: properly speaking, there is no enlightenment among demons---only speech. & \textbf{Common teaching} & None recorded \\
\addlinespace[1.5pt]
\textbf{4} & Whether the good angels have precedence over the bad angels? (\emph{SC:} Augustine (De Trinitate III) and Gregory the Great (Hom. XXXIV in Ev.)) & Affirmative: good angels have precedence over and govern the bad angels, because all authority derives from nearness to God---and beatified angels, united to God in glory, surpass demons who are estranged from Him. & \textbf{Common teaching} & None recorded \\
\end{longtable}
}

\Needspace{14\baselineskip}
\paragraph{I q.~110: How Angels Act on Bodies (Sec.~10)}
{\scriptsize
\setlength{\tabcolsep}{2.5pt}
\renewcommand{\arraystretch}{1.15}
\begin{longtable}{>{\raggedright\arraybackslash}p{0.05\linewidth}>{\raggedright\arraybackslash}p{0.23\linewidth}>{\raggedright\arraybackslash}p{0.31\linewidth}>{\raggedright\arraybackslash}p{0.12\linewidth}>{\raggedright\arraybackslash}p{0.24\linewidth}}
\toprule
\textbf{Art.} & \textbf{Question Posed \& \emph{Sed Contra}} & \textbf{Aquinas's Determination} & \textbf{Grade} & \textbf{Contrary Opinion} \\
\midrule
\endfirsthead
\toprule
\textbf{Art.} & \textbf{Question Posed \& \emph{Sed Contra}} & \textbf{Aquinas's Determination} & \textbf{Grade} & \textbf{Contrary Opinion} \\
\midrule
\endhead
\bottomrule
\endfoot
\textbf{1} & Whether the corporeal creature is governed by the angels? (\emph{SC:} Augustine (De Trinitate III.4) and Gregory the Great (Dialogues IV.6)) & Affirmative: all corporeal creatures are governed by angels. & \textbf{Common teaching} & None recorded \\
\addlinespace[1.5pt]
\textbf{2} & Whether corporeal matter obeys the mere will of an angel? (\emph{SC:} Augustine (De Trinitate III.8, cited without book title in English text)) & Negative: corporeal matter does not obey the mere nod/will of an angel so as to receive substantial forms directly from an angel. & \textbf{Common teaching} & Avicenna attributed direct formal causation in matter to separate intelligences \\
\addlinespace[1.5pt]
\textbf{3} & Whether bodies obey the angels as regards local motion? (\emph{SC:} Augustine, De Trinitate III.8--9) & Affirmative: corporeal bodies immediately obey angels as regards local motion. & \textbf{Common teaching} & None recorded \\
\addlinespace[1.5pt]
\textbf{4} & Whether angels can work miracles? (\emph{SC:} Psalm 135:4 (Vulgate)) & Negative (by their own power): God alone can work true miracles. A true miracle is an effect outside the order of the entire created nature (not merely outside one particular nature, like throwing a stone upward); everything an angel does by its own power falls within the created order. & \textbf{Common teaching} & None recorded \\
\end{longtable}
}

\Needspace{14\baselineskip}
\paragraph{I q.~111: The Action of the Angels on Man (Sec.~10)}
{\scriptsize
\setlength{\tabcolsep}{2.5pt}
\renewcommand{\arraystretch}{1.15}
\begin{longtable}{>{\raggedright\arraybackslash}p{0.05\linewidth}>{\raggedright\arraybackslash}p{0.23\linewidth}>{\raggedright\arraybackslash}p{0.31\linewidth}>{\raggedright\arraybackslash}p{0.12\linewidth}>{\raggedright\arraybackslash}p{0.24\linewidth}}
\toprule
\textbf{Art.} & \textbf{Question Posed \& \emph{Sed Contra}} & \textbf{Aquinas's Determination} & \textbf{Grade} & \textbf{Contrary Opinion} \\
\midrule
\endfirsthead
\toprule
\textbf{Art.} & \textbf{Question Posed \& \emph{Sed Contra}} & \textbf{Aquinas's Determination} & \textbf{Grade} & \textbf{Contrary Opinion} \\
\midrule
\endhead
\bottomrule
\endfoot
\textbf{1} & Whether an angel can enlighten man? (\emph{SC:} Pseudo-Dionysius, De Coelesti Hierarchia IV) & Affirmative: angels enlighten men in two ways corresponding to human cognitive nature: (1) by strengthening the human intellect's natural light; (2) by proposing intelligible truth under sensible similitudes/phantasms, since man understands only by turning to phantasms (citing Dionysius, Coel. Hier. I: ``It is impossible for the divine ray to shine on us, otherwise than shrouded by the variety of the sacred veils''). & \textbf{Common teaching} & None recorded \\
\addlinespace[1.5pt]
\textbf{2} & Whether the angels can change the will of man? (\emph{SC:} Proverbs 21:1) & Negative from within; affirmative only from without (without necessity): God alone moves the human will internally, while an angel moves it externally in two non-coercive ways---by persuading the intellect with an apprehended good and by stirring bodily passions in the sensitive appetite, leaving the will free. & \textbf{Common teaching} & None recorded \\
\addlinespace[1.5pt]
\textbf{3} & Whether an angel can change man's imagination? (\emph{SC:} Matthew 1:20; 2:13, 19 (angelic dreams of St. Joseph)) & Affirmative: both good and bad angels can move the human imagination by their natural power---not by creating brand-new sensory forms ex nihilo, but by locally moving bodily `animal spirits and humors' (citing Aristotle, De Insomniis III) so as to recombine previously received sensory impressions, whether in sleep, ecstasy, or waking life. & \textbf{Common teaching} & None recorded \\
\addlinespace[1.5pt]
\textbf{4} & Whether an angel can change the human senses? (\emph{SC:} Genesis 19:11 and 4 Kings 6:18) & Affirmative: by natural power an angel can change human external senses in two ways: (1) from without, by presenting an external sensible object (natural or angelically fashioned, like an assumed body, I q. 51 a. 2); (2) from within, by locally disturbing bodily spirits and humors (just as bile makes a sick tongue taste everything as bitter). & \textbf{Common teaching} & None recorded \\
\end{longtable}
}

\Needspace{14\baselineskip}
\paragraph{I q.~112: The Mission of the Angels (Sec.~10)}
{\scriptsize
\setlength{\tabcolsep}{2.5pt}
\renewcommand{\arraystretch}{1.15}
\begin{longtable}{>{\raggedright\arraybackslash}p{0.05\linewidth}>{\raggedright\arraybackslash}p{0.23\linewidth}>{\raggedright\arraybackslash}p{0.31\linewidth}>{\raggedright\arraybackslash}p{0.12\linewidth}>{\raggedright\arraybackslash}p{0.24\linewidth}}
\toprule
\textbf{Art.} & \textbf{Question Posed \& \emph{Sed Contra}} & \textbf{Aquinas's Determination} & \textbf{Grade} & \textbf{Contrary Opinion} \\
\midrule
\endfirsthead
\toprule
\textbf{Art.} & \textbf{Question Posed \& \emph{Sed Contra}} & \textbf{Aquinas's Determination} & \textbf{Grade} & \textbf{Contrary Opinion} \\
\midrule
\endhead
\bottomrule
\endfoot
\textbf{1} & Whether the angels are sent on works of ministry? (\emph{SC:} Exodus 23:20) & Affirmative: angels are sent by God on works of ministry. & \textbf{Church teaching} & None recorded \\
\addlinespace[1.5pt]
\textbf{2} & Whether all the angels are sent in ministry? (\emph{SC:} Gregory the Great (Hom. XXXIV in Ev.) quoting Pseudo-Dionysius (Coel. Hier. XIII)) & Negative (without exception, following Dionysius Coel. Hier. XIII): the superior angels are never sent on external ministries. & \textbf{Thomist position} & Gregory the Great (Hom. 34.12--13) allows higher angels to be sent for great mysteries while remaining within God \\
\addlinespace[1.5pt]
\textbf{3} & Whether all the angels who are sent, assist? (\emph{SC:} Gregory the Great, Moralia in Job XVII (on Job 25:3)) & Distinction: in a broad sense all good angels, even those sent on ministry, assist because all immediately gaze on God's essence (Tob 12:15); in the strict sense, only the first hierarchy assists immediately to receive divine mysteries without external mission. & \textbf{Common teaching} & None recorded \\
\addlinespace[1.5pt]
\textbf{4} & Whether all the angels of the second hierarchy are sent? (\emph{SC:} Pseudo-Dionysius, De Coelesti Hierarchia VIII) & Negative: not all angels of the second hierarchy are sent externally. & \textbf{Thomist position} & Gregory places Virtues in the lowest hierarchy (Hom. 34.7) \\
\end{longtable}
}

\Needspace{14\baselineskip}
\paragraph{I q.~113: The Guardianship of the Good Angels (Sec.~10)}
{\scriptsize
\setlength{\tabcolsep}{2.5pt}
\renewcommand{\arraystretch}{1.15}
\begin{longtable}{>{\raggedright\arraybackslash}p{0.05\linewidth}>{\raggedright\arraybackslash}p{0.23\linewidth}>{\raggedright\arraybackslash}p{0.31\linewidth}>{\raggedright\arraybackslash}p{0.12\linewidth}>{\raggedright\arraybackslash}p{0.24\linewidth}}
\toprule
\textbf{Art.} & \textbf{Question Posed \& \emph{Sed Contra}} & \textbf{Aquinas's Determination} & \textbf{Grade} & \textbf{Contrary Opinion} \\
\midrule
\endfirsthead
\toprule
\textbf{Art.} & \textbf{Question Posed \& \emph{Sed Contra}} & \textbf{Aquinas's Determination} & \textbf{Grade} & \textbf{Contrary Opinion} \\
\midrule
\endhead
\bottomrule
\endfoot
\textbf{1} & Whether men are guarded by the angels? (\emph{SC:} Psalm 90:11 (Vulgate)) & Affirmative: men are guarded by angels. & \textbf{Church teaching} & None recorded \\
\addlinespace[1.5pt]
\textbf{2} & Whether each man is guarded by an angel? (\emph{SC:} Jerome (Commentary on Matthew 18:10)) & Affirmative: each individual human being has a personal guardian angel appointed to him. & \textbf{Church teaching} & None recorded \\
\addlinespace[1.5pt]
\textbf{3} & Whether to guard men belongs only to the lowest order of angels? (\emph{SC:} Psalm 90:11 with Pseudo-Dionysius (Coel. Hier. V, IX)) & Distinction: particular guardianship (one angel assigned to each individual man) belongs strictly to the lowest order (`Angels,' who announce lesser things, Gregory Hom. XXXIV). & \textbf{Thomist position} & Chrysostom (In Matt. Hom. 59.4) speaks of greater angels guarding believers without fixing a single choir \\
\addlinespace[1.5pt]
\textbf{4} & Whether angels are appointed to the guardianship of all men? (\emph{SC:} Jerome (Commentary on Matthew 18:10, cited via a. 2)) & Affirmative for all earthly wayfarers (except Christ): as long as any human being journeys along the perilous road of earthly life (Ps 141:4), he has a guardian angel---including Adam in innocence, unbelievers, the reprobate, and even Antichrist. & \textbf{Common teaching} & None recorded \\
\addlinespace[1.5pt]
\textbf{5} & Whether an angel is appointed to guard a man from his birth? (\emph{SC:} Jerome (Commentary on Matthew 18:10)) & Affirmative: siding with Jerome over the alternative baptism-start opinion noted by Origen (Tract. V super Matt.), Aquinas holds that a personal guardian angel is appointed at birth, because angelic guardianship is a benefit bestowed on man as a rational creature (not solely as a Christian, unlike the sacraments beginning at baptism). & \textbf{Disputed} & Origen (In Matt. XIII.26--29) leaves open whether personal guardianship begins at birth or at baptism \\
\addlinespace[1.5pt]
\textbf{6} & Whether the angel guardian ever forsakes a man? (\emph{SC:} 1 Peter 5:8 (a fortiori from demonic assault)) & Negative as to total abandonment; affirmative only in particular permissions: a guardian angel never wholly forsakes a man (just as God's providence never wholly withdraws from any being), though in accordance with divine judgments he sometimes permits a man to suffer tribulation or fall into sin (explaining Jer 51:9 on Babylon and Isa 5:5 + Gloss). & \textbf{Common teaching} & None recorded \\
\addlinespace[1.5pt]
\textbf{7} & Whether angels grieve for the ills of those whom they guard? (\emph{SC:} Apocalypse 21:4) & Negative: angels do not grieve over either the sins or the punishments of those they guard. & \textbf{Common teaching} & None recorded \\
\addlinespace[1.5pt]
\textbf{8} & Whether there can be strife or discord among the angels? (\emph{SC:} Daniel 10:13) & Negative as to opposition of wills (all good angels share perfect peace and concord in God's decree, Job 25:2); affirmative only as to intellectual consultation over competing merits. & \textbf{Common teaching} & None recorded \\
\end{longtable}
}

\Needspace{14\baselineskip}
\paragraph{I q.~114: The Assaults of the Demons (Sec.~10)}
{\scriptsize
\setlength{\tabcolsep}{2.5pt}
\renewcommand{\arraystretch}{1.15}
\begin{longtable}{>{\raggedright\arraybackslash}p{0.05\linewidth}>{\raggedright\arraybackslash}p{0.23\linewidth}>{\raggedright\arraybackslash}p{0.31\linewidth}>{\raggedright\arraybackslash}p{0.12\linewidth}>{\raggedright\arraybackslash}p{0.24\linewidth}}
\toprule
\textbf{Art.} & \textbf{Question Posed \& \emph{Sed Contra}} & \textbf{Aquinas's Determination} & \textbf{Grade} & \textbf{Contrary Opinion} \\
\midrule
\endfirsthead
\toprule
\textbf{Art.} & \textbf{Question Posed \& \emph{Sed Contra}} & \textbf{Aquinas's Determination} & \textbf{Grade} & \textbf{Contrary Opinion} \\
\midrule
\endhead
\bottomrule
\endfoot
\textbf{1} & Whether men are assailed by the demons? (\emph{SC:} Ephesians 6:12) & Affirmative: men are assailed by demons. & \textbf{Church teaching} & None recorded \\
\addlinespace[1.5pt]
\textbf{2} & Whether to tempt is proper to the devil? (\emph{SC:} 1 Thessalonians 3:5 (with Gloss)) & Affirmative (in the specific sense of tempting in order to harm and lead into sin): properly, to tempt is to test so as to gain knowledge, directed to a good or bad end. & \textbf{Common teaching} & None recorded \\
\addlinespace[1.5pt]
\textbf{3} & Whether all sins are due to the temptation of the devil? (\emph{SC:} Gennadius / Pseudo-Augustine, De Ecclesiasticis Dogmatibus XLIX) & Distinction: indirectly and occasionally, the devil is the cause of all human sins by instigating the Fall of our first parents, which introduced universal proneness to sin (saving Dionysius, Div. Nom. IV, and Damascene, De Fide Orth. II.4). & \textbf{Common teaching} & None recorded \\
\addlinespace[1.5pt]
\textbf{4} & Whether demons can lead men astray by means of real miracles? (\emph{SC:} Pseudo-Augustine (Liber XXI Sententiarum, sent. 4, cited as Augustine Q. 83)) & Negative in the strict sense of miracle (God alone acts outside the whole created order, I q. 110 a. 4); affirmative in the wide sense of marvels exceeding human power and experience---some of which are real physical effects produced via natural seeds (Pharaoh's magicians' serpents/frogs; fire and storm against Job, Augustine De Civ. Dei XX.19), while transformations exceeding nature (turning a man into a beast, raising the dead) are mere sensory/imaginative illusions. & \textbf{Common teaching} & None recorded \\
\addlinespace[1.5pt]
\textbf{5} & Whether a demon who is overcome by man, is for this reason hindered from making further assaults? (\emph{SC:} Matthew 4:11) & Qualified affirmative (temporarily, not permanently): siding with the view that a defeated demon withdraws from the same person for a definite time (Luke 4:13: ``departed from Him for a time''), Aquinas gives two reasons: (1) God's clemency limiting temptation to human strength (Opus Imperfectum in Matt. Hom. V, cited as Chrysostom); (2) demonic astuteness shrinking from repeated defeat (Ambrose on Luke 4:13)---yet the demon can return later (Matt 12:44). & \textbf{Common teaching} & None recorded \\
\end{longtable}
}
